% TOP_PROBLEM: {"id": 30002380, "title": "Serre's Uniformity Problem for Elliptic Curves", "category_id": 1, "category": {"id": 1, "name": "number_theory", "display_name": "Number Theory", "description": "Properties of integers, prime numbers, Diophantine equations.", "slug": "number-theory", "order_index": 1, "created_at": "2026-07-31T15:26:25.670Z"}, "status": "partially_solved"}
% FIRST_POSTED: 2026-10-01
% ATTEMPT_STATUS: unresolved
% !TeX program = lualatex
% Definition and sources only; this file is not an LLM solution attempt.
\documentclass[11pt]{article}
\usepackage[margin=25mm]{geometry}
\usepackage{fontspec}
\setmainfont{Latin Modern Roman}
\usepackage{amsmath,amssymb,mathtools}
\usepackage{unicode-math}
\setmathfont{Latin Modern Math}
\usepackage{xurl,hyperref}
\hypersetup{colorlinks=true,urlcolor=blue}
\setcounter{secnumdepth}{0}
\setlength{\parindent}{0pt}
\setlength{\parskip}{5pt}
\setlength{\emergencystretch}{3em}
\begin{document}
\section*{\texorpdfstring{Serre's Uniformity Problem for Elliptic Curves}{Serre's Uniformity Problem for Elliptic Curves}}\label{30002380}
\subsection{Definitions and mathematical statement}
For an elliptic curve $E/{\mathbb{Q}}$ and a prime $p$, its geometric $p$-torsion is a two-dimensional ${\mathbb{F}}_p$-vector space. Choosing a basis gives
$\rho_{E,p}:\operatorname{Gal}(\overline{{\mathbb{Q}}}/{\mathbb{Q}})\to\mathrm{GL}_2({\mathbb{F}}_p)$, well defined up to conjugation. Non-CM means $\operatorname{End}_{\overline{{\mathbb{Q}}}}(E)={\mathbb{Z}}$. The question is
\[
 \exists N>0\quad\forall E/{\mathbb{Q}}\text{ non-CM}\quad\forall p>N,
 \qquad\operatorname{im}\rho_{E,p}=\mathrm{GL}_2({\mathbb{F}}_p).
\]
The bound must be independent of the curve, not merely finite for each fixed curve. It is surjectivity of the full residual representation, not just irreducibility or a condition on its determinant.
\par\smallskip\textit{Formulation and concept references:} \href{https://arxiv.org/abs/2305.17780}{[S1]}.

\subsection{Short English statement}
Is there a single prime cutoff beyond which every non-CM elliptic curve over the rationals has the largest possible Galois action on its torsion points?
\subsection{Research attempt: a transvection criterion and the remaining arithmetic obstruction}
\textbf{Status.} We prove a finite-group surjectivity criterion and an explicit counterexample to weakening it to irreducibility. This isolates an arithmetic condition that would suffice for uniformity; it does not prove that condition uniformly in elliptic curves. The source [S1] analyzes the residual nonsplit-Cartan-normalizer case and must not be read as an unconditional solution of the selected quantifier statement.

\subsubsection{1. Irreducibility plus a transvection forces the special linear group}
Let $H\le\mathrm{GL}_2(\mathbb F_p)$ act irreducibly and contain a nonidentity transvection $T$, meaning $(T-I)^2=0$ and $T\ne I$. Its fixed space is a line $L$. Irreducibility supplies $h\in H$ with $hL\ne L$; otherwise $L$ would be invariant. The conjugate $hTh^{-1}$ fixes the distinct line $M=hL$.

In a basis whose first vector spans $L$ and whose second spans $M$, the two matrices have forms
\[U(a)=\begin{pmatrix}1&a\\0&1\end{pmatrix},\qquad
V(b)=\begin{pmatrix}1&0\\b&1\end{pmatrix},\quad a,b\ne0.\]
Their powers give every $U(t)$ and $V(t)$, $t\in\mathbb F_p$, because the additive group of the prime field is generated by each nonzero element. Elementary row operations using these matrices reduce any determinant-one matrix to a diagonal matrix. The diagonals are also generated: with
\[w(t)=U(t)V(-t^{-1})U(t)=\begin{pmatrix}0&t\\-t^{-1}&0\end{pmatrix},\]
we have $w(t)w(-1)=\operatorname{diag}(t,t^{-1})$. Hence $\mathrm{SL}_2(\mathbb F_p)\subseteq H$. This proof includes $p=2$ and does not appeal to a classification of all subgroups.

If additionally $\det(H)=\mathbb F_p^\times$, then $H=\mathrm{GL}_2(\mathbb F_p)$. For any $A$ choose $h\in H$ with $\det h=\det A$; then $Ah^{-1}\in\mathrm{SL}_2\subseteq H$. For an elliptic curve over $\mathbb Q$, the Weil pairing identifies the determinant with the cyclotomic character, whose mod-$p$ image is all of $\mathbb F_p^\times$. Thus irreducibility and a transvection suffice for the actual representation.

\subsubsection{2. Why irreducibility and the determinant alone do not suffice}
Regard $\mathbb F_{p^2}$ as a two-dimensional vector space over $\mathbb F_p$. Multiplication by nonzero elements embeds the group $C=\mathbb F_{p^2}^\times$ into $\mathrm{GL}_2(\mathbb F_p)$. If a one-dimensional subspace were stable under every multiplication, then for any nonzero vector $v$ in it, the set $\mathbb F_{p^2}^\times v$ would stay in that line, although it consists of every nonzero vector. Thus the action is irreducible.

Its determinant is the norm $a\mapsto a^{p+1}$, surjective onto $\mathbb F_p^\times$ because the multiplicative group of a finite field is cyclic. Nevertheless
\[|C|=p^2-1<|\mathrm{GL}_2(\mathbb F_p)|=(p^2-1)p(p-1).\]
This subgroup has no element of order $p$, so no nontrivial transvection. It provides a concrete algebraic obstruction to any proof that checks only irreducibility and determinant. It is not asserted to be the Galois image of a non-CM rational elliptic curve for arbitrarily large $p$.

\subsubsection{3. The uniformity gap}
A sufficient arithmetic theorem would give one bound $N$ such that every non-CM $E/\mathbb Q$ and every $p>N$ has irreducible image containing a transvection. Even if this criterion is stronger than necessary, the group proof shows exactly how it would conclude surjectivity. Available curve-dependent estimates cannot be inserted here by choosing $N=N(E)$: the target quantifies $N$ before $E$. Likewise, excluding one proper subgroup family leaves the remaining families to be treated.

\subsubsection{4. Adversarial verification}
The transvection argument requires the prime field: over a larger finite field, powers of one unipotent element need not fill its entire root subgroup. It also needs an actual element of order $p$, not merely a matrix with determinant one. The diagonal formula and group orders were checked for small primes in the accompanying diagnostic. No arithmetic argument in this attempt rules out the Cartan alternatives uniformly, so the requested global cutoff remains unproved.

\subsection{Sources}
\begin{itemize}
\item[S1]Formal statement, abstract.
\url{https://arxiv.org/abs/2305.17780}.
\textit{Formulation source.}
\end{itemize}
\end{document}
